\section{Introduction}
It had been argued that certain axio-dilaton couplings in the 4-point, 10d, type IIB effective action may be viewed from a 12d perspective.
In \cite{Minasian:2015bxa}, it was suggested that the 12d $\hat{t}_8 \hat{t}_8 \hat{R}^4$ exactly produces the 4-point axio-dilaton couplings at 8 derivative.
However, such study is in fact somewhat limited.



step 1 and 2 are correct, but step 3 is a little ambitious. Right now we don't even know if there are enough terms to be packaged into a single 12d quartic invariant.

For example, let us consider an object like

(R_{MNRS}R^{MNRS})^2

upon reduction, taking into account only the nonvanishing components of R_{MNRS}, we have

(R_{mnrs}R^{mnrs} + 2 R_{ambn}R^{ambn} + 2 R_{mnab}R^{mnab} + R_{abcd}R^{abcd})^2


We know that (R_{mnrs}R^{mnrs})^2 exists, and I myself have found 2R_{mnrs}R^{mnrs}R_{ambn}R^{ambn}, the rest are unknown. Finding a single such term would be a very nice find.

However, there are difficulties, for a term like (R_{MNRS}R^{MNRS})^2, because certain components of R_{ambn} can be quadratic in axiodilaton, this term will contribute all the way to 8-point type IIB effective actions, so at 4-point, we are not examining all the effective actions out there, and 5-point perhaps better, but maybe not that much, because beginning at 5-point we have observed terms that can not be uplifted to 12d Riemann tensor: the PP terms. Though it is posssible that there are further terms that can not be uplifted to 12d Riemann tensor, but when combined they do get uplifted.
